\chapter{Projektplan}

\includegraphics{appendix/img/openhabLogo}

\begin{description}[style=multiline,leftmargin=3cm]
\item[Thema] Aufbau einer Smart-Home Beispielapplikation
\item[Studenten] Dominik Freier, Marco Leutenegger
\item[Betreuer] Prof. Hansjörg Huser
\end{description}
\pagebreak

\section*{Änderungsgeschichte}
	\begin{tabularx}{\textwidth}{lllX}
	\textbf{Datum}		& \textbf{Version}	& \textbf{Änderung}	& \textbf{Autor} \\
	\hline
	25.02.2015			& 0.0.1				& Dokument erstellen & Marco Leutenegger \\
	\hline
	\tbd					& 0.0.2				& Dokument aktualisieren	& Dominik Freier, Marco Leutenegger \\
	\hline
	\end{tabularx}
\pagebreak

\section*{Einführung}
\subsection*{Zweck}
Dieses Dokument dient als Projektplan für die Bachelorarbeit von Dominik Freier und Marco Leutenegger und definiert alle organisatorischen Rahmenbedingungen.

\subsection*{Gültigkeitsbereich}
Die Gültigkeit des Projektplans beschränkt sich auf die Bachelorarbeit von Dominik Freier und Marco Leutenegger im Frühjahrssemester 2015.

\subsection*{Referenzen}
\begin{tabularx}{\textwidth}{lX}
	\textbf{Bezeichnung}	& \textbf{Referenz} \\
	\hline
	Risikomanagement		& Siehe separates Dokument\\
	\hline
	Security Infos		& \url{https://github.com/openhab/openhab/wiki/Security} \\
	\hline
\end{tabularx}
\pagebreak

\section*{Projekt und Übersicht}
\subsection*{Zweck und Ziel}
Diese Bachelorarbeit hat as Ziel, eine Smart-Home Beispielapplikation aufzubauen, welche wesentliche Aspekte einer Internet-of-Things-Anwendung demonstriert, wie Steuern von Devices, Lesen von Sensoren, Event-Verarbeitung, Überwachung und intelligente Abläufe steuern, Streaming von Sensordaten und Online-Analyse der Daten usw. \\
Das System soll auf einer tragfähigen und erweiterbaren Architektur aufgebaut werden und Microsoft Azure als Cloud Plattform benutzen.

\subsection*{Lieferumfang}
Die abzuliefernden Dokumente und Software-Artefakte des Projekts richten sich im Wesentlichen nach den Vorgaben aus den Dokumentationsanleitungen der HSR. Eine davon abweichender Lieferumfang wurde mit dem Betreuer besprochen und genehmigt.

\subsection*{Referenzen}
\begin{tabularx}{\textwidth}{llXll}
	\textbf{Nr.}	& \textbf{Art} & \textbf{Bezeichnung} & \textbf{Form} & \textbf{Empfänger}\\
	\hline
	1 & Publikation & Poster 							& PDF 			& H.Huser \\	\hline
	2 & Publikation & Kurzfassung 						& PDF 			& H.Huser \\\hline
	3 & Dokument		& Bericht 							& PDF/Ausdruck	& H.Huser \\\hline
	4 & Dokument		& Projektplan 						& PDF/Ausdruck	& H.Huser \\\hline
	5 & Dokument		& Sitzungsprotokolle 				& PDF/Ausdruck	& H.Huser \\\hline
	6 & Dokument		& Eigenständigkeitserklärung			& PDF/Ausdruck	& H.Huser \\\hline
	7 & Dokument		& Erfahrungsbericht D.Freier 		& PDF/Ausdruck	& H.Huser \\\hline
	8 & Dokument		& Erfahrungsbericht M.Leutenegger 	& PDF/Ausdruck	& H.Huser \\\hline
	9 & Source		& Code-Abgabe 						& ZIP			& H.Huser \\\hline
	10& Archiv 		& 2x Deliverables 1-9 				& DVD			& H.Huser \\\hline
\end{tabularx}
\pagebreak


\section*{Projektorganisation}
Die Dokumentation des Projekts gliedert sich in diesen Projektplan und einen Bericht. Im Projektplan werden alle organisatorischen Aspekte festgehalten, wie etwa die Planung der Meilensteine, Aufgaben der Teammitglieder oder Abmachungen zum Dokumentemanagement. Im Bericht werden technische Beschreibungen der Ausgangslage, Diskussionen für Lösungsansätze, Requirements und Details zur Umsetzung dokumentiert. \\
\\
Damit die Teammitglieder möglichst parallel und effizient arbeiten können, werden alle Dokumente mit LaTeX geschrieben und auf einem Git-Repository verwaltet. Daruch wird das Risiko von Versionskonflikten reduziert und der Zugriff insbesondere für den Betreuer vereinfacht. \\
\\
Die Verwaltung der Aufgaben und agilen Vorgänge erfolgt durch Jira. Wir erhielten zu diesem Zweck eine Classroom Lizenz vom Hersteller Atlassian. Jira wurde auf einem virtuellen Server der HSR installiert.

\subsection*{Organisationsstruktur}
\begin{tabularx}{\textwidth}{Xl}
	\textbf{Verantwortung}	& \textbf{Teammitglied} \\
	\hline
	Verwaltung und Bereinigung de Dokumente		& D. Freier, M. Leutenegger \\ \hline
	Pflege virtueller Server, Jira inkl. Backups	& M. Leutenegger \\ \hline
	Sitzungsprotokolle verfassen					& D. Freier, M. Leutenegger \\ \hline
	Iterationsplanung							& D. Freier, M. Leutenegger \\ \hline
	Risikomanagement								& D. Freier, M. Leutenegger \\ \hline
	Architekturdesign							& D. Freier, M. Leutenegger \\ \hline
\end{tabularx}

\subsection*{Externe Schnittstellen}
Betreuer der Bachelorarbeit ist Prof. Hansjörg Huser. Experte ist Herr Stefan Zettel. Gegenleser ist \tbd .
\pagebreak

\section*{Management Abläufe}
\subsection*{Zeitliche Planung}
Das Projekt wird während des Frühjahrssemester 2015 durchgeführt. Der Start der Arbeit war am Montag, den 16. Februar 2015. Die Abgabe der Vollständigen Dokumentation an den Betreuer erfolgt am Freitag, den 12. Juni 2015. Als Zeitbudget sollen in den 17 Wochen insgesamt 720 Stunden, bzw. rund 21 Stunden pro Woche und Student eingeplant werden.

\subsection*{Vorgehensmodell}
Als Vorgehensmodell wurde der Rational Unified Process ausgewählt, da das Projektteam mit diesem Modell aus früheren Arbeiten (inkl. Semesterarbeit) vertraut ist und damit gute Erfahrungen gemacht hat. Die Phasen wurden nach dem Schema "eins, drei, zwei, ein" in insgesamt sieben Iterationen à zwei Wochen aufgeteilt. \tbd

\subsection*{Meilensteine}
\tbd

\subsection*{Iterationsplanung}
\tbd

\subsection*{Besprechungen}
Wöchentliche Besprechungen: \\

\begin{tabularx}{\textwidth}{lXlll}
	\textbf{Bezeichnung}	& \textbf{Ziel} & \textbf{Wochentag} & \textbf{Uhrzeit} & \textbf{Ort}\\
	\hline
	Teambesprechung			& Projektarbeiten im Plenum erledigen	& Donnerstag & 08:10-08:40	& HSR (Labor) \\ \hline
	Fortschrittsbesprechung	& Fortschritte / Probleme besprechen		& Mittwoch & 10:10-10:50	& HSR (6.010) \\ \hline
\end{tabularx}
\pagebreak

\section*{Risikomanagement}
\subsection*{Risiken}
Unser projektbezogenes Risikomanagement wird in einem separaten Excel-Sheet aufgeführt. Nachstehend ist ein Snapshot dieses Sheets zu sehen. Das Dokument wird laufend angepasst und aktualisiert, falls notwendig. \tbd

\subsection*{Umgang mit Risiken}
\subsubsection*{Reserven / Rückstellungen}
Das grösste Risiko stellt R3 (ungeplante Machbarkeiten) \tbd dar. Aus diesem Grund werden in diesem Projekt Rückstellungen von 20 Stunden eingeplant.

\subsubsection*{Überprüfung von Risiken}
Weitere Risiken werden im Laufe des Entwicklungsprozesses erkennbar. Hierfür verwenden wir ein Entscheidungs-Dokument, welches als zentrale Stelle dient, um Entscheidungen und Risiken zu DOkumentieren und eine zentrale Anlaufstelle bei Fragen darstellt. Des weiteren wird in der Beschreibung des betroffenen Vorgangs auf mögliche Risiken hingewiesen und dokumentiert.
\pagebreak

\section*{Arbeitspakete}
Die Arbeitspakete wurden im Projektmanagementtool Jira als Vorgänge definiert. \\
Einige Vorgänge beinhalten weitere Untertätigkeiten, die wir ebenfalls als einzelne Arbeitspakete betrachten. \\
\\
Eine Übersicht mit allen Arbeitspaketen und dem zeitlichen Ablauf nach Iterationen befindet sich unter: \url{http://sinv-56046.edu.hsr.ch:8080} > Agile > Zeige alle Boards > baIOTBoard > Plan

\section*{Infrastruktur}
\subsection*{Software}
Wie in jedem Projekt kommt verschiedene Software zum Einsatz. \\

\begin{tabularx}{\textwidth}{lXlll}
	\textbf{Bezeichnung}	& \textbf{Ziel} & \textbf{Wochentag} & \textbf{Uhrzeit} & \textbf{Ort}\\
	\hline
	Teambesprechung			& Projektarbeiten im Plenum erledigen	& Donnerstag & 08:10-08:40	& HSR (Labor) \\ \hline
	Fortschrittsbesprechung	& Fortschritte / Probleme besprechen		& Mittwoch & 10:10-10:50	& HSR (6.010) \\ \hline
\end{tabularx}

\section*{Qualitätsmassnahmen}
\tbd

\subsection*{Dokumentation}
\subsubsection*{Ablage}
\tbd

\subsubsection*{Qualität}
\tbd

\subsection*{Projektmanagement}
\tbd

\subsection*{Entwicklung}
\subsubsection*{Code Reviews}
\tbd

\subsubsection*{Code Style Guidelines}